% Canonical editable manuscript. Build with experiments/paper/build_sigir.py.
\makeatletter
\def\input@path{{template/}}
\makeatother
\documentclass[sigconf,natbib=true,anonymous=true,pbalance=true]{acmart}
\usepackage{booktabs}
\usepackage{amsmath}
\usepackage{xurl}
\usepackage{array}
\emergencystretch=2em
\setcopyright{none}
\settopmatter{printacmref=false}
\acmConference[Research draft]{Research draft}{November 2026}{}
\acmYear{2026}
\renewcommand\footnotetextcopyrightpermission[1]{}
\title[Low-Cost Evidence-Based Expert Search]{Low-Cost First-Query Expert Search with Attributed Evidence}
\author{Anonymous Authors}
\begin{document}

\begin{abstract}
Expert search must find people whose own contributions meet a request, often with support in different publications. We study a first-query agent that joins ranked source pools by contributor identity and uses indexed personal roles to select a small set of declarations for verification. In a 452-article mixed PLOS and eLife catalog, at most three reads per request recover 85 of 160 PLOS and 46 of 51 eLife label-supported matches at candidate depth 128, compared with 45 and 14 for byline-only joint nomination. Raising depth from 64 to 128 increases coverage, while the joint role route returns the same experts and citations on requests answered at both depths. A shallow-planning certificate preserves the fixed-depth responses and reduces reads from 633 to 516 with complete role fields. An independent byline precheck reduces reads from 1,266 to 910 in a paired live retrieval and source-service replay with identical responses. The source service uses loopback HTTP and an inserted read delay. These studies support fixed-budget coverage and equal-output execution savings. Independent evidence judgments and free-form enterprise requests remain necessary to establish expert-ranking quality.
\end{abstract}
\ccsdesc[500]{Information systems~Information retrieval}
\ccsdesc[300]{Information systems~Retrieval models and ranking}
\keywords{expert search, evidence verification, attribution, early stopping, evaluation}
\maketitle

\section{Introduction}
An enterprise looking for help with industrial inspection may need a person who has worked under a specified line condition and can document a measured operating result. A paper about the same equipment supplies a retrieval lead. A source stating a design target or citing someone else's experiment supplies different evidence. Expert search must retain those distinctions while finding useful candidates within a practical time budget.

Reading the individually highest-ranked sources can spend most of the budget on papers about different people. A last read may then be unable to complete either person's evidence. We move the same-person constraint into the initial source plan. Ranked topic lists and contributor identifiers nominate complementary publications. The source declaration must still establish that the candidate performed each requested duty.

Our first contribution is contributor-conditioned source allocation at a fixed verification budget. We measure reachability, loaded evidence and the final citation choice separately, and compare both natural and matched issuance. Our second contribution applies exact pruning to this allocation. A byline precheck removes a planned read that cannot contribute to any proof, and a prefix certificate can defer deeper role planning without changing the fixed-depth outcome. The evaluation separates CPU ranking, local declaration access and a controlled HTTP source service. We also examine attribution errors and controlled prose interpretation. Continual learning remains a follow-up to the frozen first-query agent.

\section{Related Work}
Classical expert search aggregates document relevance into a ranking of people~\cite{balog2006formal}. We retain that aggregation for candidate generation and use contributor declarations to guide nomination and verification. CoF evaluates reviewer--paper relevance~\cite{zhang2025cof}. MERIT distinguishes core expertise from incidental overlap~\cite{yang2026merit}. Neither signal supplies ground truth for a person's contribution to a cited work in our enterprise setting.

Rank joins exploit input ordering to produce high-scoring joins~\cite{ilyas2003rankjoin}. Reciprocal-rank fusion combines the topical ranks from lexical and dense retrieval~\cite{cormack2009rrf}. We use a specialized join to allocate reads before retrieving declarations. SetR selects passages that jointly meet information requirements~\cite{lee2025setr}. Our comparison imposes an identity constraint across distinct sources and verifies support at the expert level. RATE learns reviewer profiles for expertise ranking~\cite{liu2026rate}. It provides a separate comparison for reviewer fit. S2G-RAG predicts evidence gaps and translates them into follow-up queries~\cite{li2026s2g}. We test whether using contributor links before the first read improves what the same source budget can recover.

Threshold algorithms use monotone bounds to avoid access when unseen values cannot change a top-$K$ answer~\cite{fagin2001aggregation,broder2003wand}. Our fixed-inventory variant also requires stability of the supporting citations. Adaptive-RAG routes retrieval by predicted query complexity~\cite{jeong2024adaptive}, and ranking cascades allocate expensive scoring to promising items~\cite{wang2011cascade,hofstatter2021intradocument}. We keep those concerns distinct from the source allocation comparison.

LitSearch provides author-written scientific search questions~\cite{ajith2024litsearch}. CSFCube provides expert grades within candidate pools for specified paper facets~\cite{mysore2021csfcube}. We use them to distinguish source acceptance from expert ranking. Source-grounded synthesis and citation evaluation distinguish coverage from citation correctness~\cite{lin2025scirgen,gao2023alce,xu2025citeeval}. TREC Enterprise already judges support at the topic--expert--document level and credits a relevant expert if any returned source supports the match~\cite{soboroff2006enterprise}. Our citation intervention keeps that criterion as a comparator and tests whether the full set of citations supports the claims about the nominated person.

\section{Task and method}
\label{sec:method}
A request $q$ identifies a target operation, context, and any hard requirements. The initial agent retrieves owned source records using a frozen index and aggregates them into a deterministic expert prior. The initial route needs no user feedback. It keeps candidates with incomplete declarations eligible for reading (Figure~\ref{fig:plan}). For each expert $e$ and requirement $r$, the index supplies an ordered list $W_{er}$ of candidate passages. The evidence state stores confirmed support, rejection, unread evidence, or an observed result that still needs rechecking. An unavailable source cannot be a returned witness. We preserve query spans, source hashes, ownership and passage offsets so that the returned decision can be audited.

\begin{figure*}[t]
\centering
\setlength{\fboxsep}{5pt}
\newcommand{\planstage}[2]{\fbox{\begin{minipage}[c][.58in][c]{.145\textwidth}\centering
\small\textbf{#1}\par\scriptsize #2\end{minipage}}}
\planstage{Requirements}{Topics and roles\\Preserve raw constraints}
\hfill\(\rightarrow\)\hfill
\planstage{Ranked pools}{Lexical and dense ranks\\One list per condition}
\hfill\(\rightarrow\)\hfill
\planstage{Joint plan}{Same contributor\\Distinct sources}
\hfill\(\rightarrow\)\hfill
\planstage{Source reads}{Precheck or certify\\Read original declarations}
\hfill\(\rightarrow\)\hfill
\planstage{Verified result}{Expert and citations\\Unresolved obligations}
\caption{The frozen first-query path. Role postings guide nomination. Original declarations establish attributed support. Source prechecks and stopping certificates preserve the specified reference outcome.}
\label{fig:plan}
\end{figure*}

\subsection{Joint nomination before reading}
\label{sec:nomination}
For each of two topic--role conditions, hybrid retrieval ranks title and abstract records. Let $D_1,D_2$ denote these lists, $r_i(d)$ a source's rank on list $i$, and $B(d)$ its indexed contributor identifiers. An eligible ordered pair $(a,b)$ has $a\in D_1$, $b\in D_2$, $a\ne b$, and $B(a)\cap B(b)\ne\varnothing$. The first two reads minimize the key
\[
\left(\max\{r_1(a),r_2(b)\},\ r_1(a)+r_2(b),\ a,\ b\right)
\]
over eligible pairs, using stable identifiers for ties. To find that pair, retain the first two sources per contributor and list. For one contributor, distinct first choices give the minimum. If both are the same source, any distinct pair must replace at least one first choice. Replacing just that choice with its second source is at least as good as any later substitution. Comparing these minima across contributors gives the scan solution.

A third read maximizes the number of additional directed eligible pairs with the first two sources. Its ties favor better topical rank, then a stable identifier. When no eligible pair exists, use the top distinct topical sources. Only the resulting sources are opened. A recommendation requires the original declarations to support both requested roles for the same verified person. Publisher topic tags are reserved for evaluation. The exact solution applies to two lists. Incomplete author metadata can reduce coverage.

\subsection{Indexed roles and exact source execution}
An author--source role posting is compatible with a condition if it lists the requested role or has unavailable role information. For an indexed author with a nonempty known role set, a different role excludes that match during nomination. The role-aware route minimizes the number of unknown roles before applying the topical pair key. For each contributor and list, the best two compatible rows in this order suffice for the distinct-source constraint. The third source maximizes fully indexed pairs with the anchors, then possible pairs. Uncertain or incorrect roles in the index can alter what gets read. Verification still tests the original declarations.

\paragraph{Prune before parallel reads.}
For the fixed three-source byline plan, an unopened third source can only contribute a proof if it can pair with one of the two other sources. Skip it when it has no common indexed author on a compatible orientation of the two ranked lists. Parallelism is retained. This precheck assumes complete author identifiers for the reference scope and verifies every other proposed witness.

\paragraph{Certify shallow role planning.}
Let the shallow lists be prefixes of the fixed deep lists. The selected shallow pair remains the deep anchor pair when both roles are indexed. It already attains the minimum number of unknown roles, and every pair involving a position outside a required shallow prefix has a worse minimax rank. Read the two anchors and consider all possible completions through one extra shallow source using its byline. Grant every shared contributor the needed role optimistically. If a verified prefix proof outranks every such completion, and therefore also every proof using a rank beyond the prefix, return it. Otherwise finish deep planning and reuse any already-open declaration. This certificate preserves the selected expert and both citations for the same indexed reference. It saves role processing and reads while the benchmark computes the full deep ranking for both routes.

\subsection{Request-level interpretation and residual proof}
A structured request specifies topical content and two contribution duties independently. Interpret each duty against the available declaration fields before binding it to a source. We run two answer encodings and accept a field only on agreement. A single GPU batch executes the four prompts for a two-duty query. The resulting fields are reused across candidate sources. Every binding still uses the particular report and its author declaration. Field comparison can replace repeated interpretation when the relevant declarations are complete.

For free text, a role proposal can influence nomination while the complete, unverified requirement spans remain available to the verifier. The system needs evidence before it can discharge a requirement, including one that appears in the same sentence as a role. Exact span offsets establish quotation identity. Conditions such as operating scale, test date and exclusive responsibility need separate support, even when the role is well understood. Our structured timing excludes free-text proof. A separate controlled prose replay encodes two topic spans and two role phrases in one forward through a frozen 0.6B encoder, with a ridge role head trained on 84 task descriptions absent from the held-out phrase inventory by exact match. The templates fix span boundaries. We compare separate forwards and open 4B interpretation while holding the downstream evidence route fixed.

\subsection{Certificates for ranked citations}
For the hard-constraint variant, an expert is confirmed when each required $r$ has at least one positive witness. A separate configuration confirms an expert after a fixed number of distinct positive sources. In both cases the prior orders confirmed experts, and required witness citations follow the published order in $W_{er}$. For the reviewer configuration, unconfirmed positions can be filled by the same fixed prior with empty citations. The hard-constraint configuration returns confirmed candidates only.

Let $z$ denote the vector of trusted binary witness decisions and $U$ the unread or tentative decisions. To obtain $R^{-}$, assign failure to every member of $U$. For $R^{+}$ assign support to every member of $U$. Both outputs use the same ranking and citation conventions. Under these rules,
\begin{equation}
R^{-}=R^{+}\quad\Longrightarrow\quad R(z,u)=R^{-}
\qquad\text{for every admissible }u.
\label{eq:certificate}
\end{equation}
For the proof, promotion is monotone in witness support. If a potentially promoted expert could displace a returned expert in any completion, that displacement would also change the confirmed order or citation status in the upper completion. The first positive witness for a requirement can only move earlier as support is added. The first $c$ positive sources for a fixed support count have the same componentwise property. Equality of the two endpoint witness lists fixes each selected citation in between. The argument uses a fixed candidate prior and monotone confirmation. It does not cover a score that penalizes a candidate when a new conflicting source appears.

At each step, the agent can read the first unseen source for an expert that remains eligible for promotion. It can also resolve an ambiguity that affects an optimistic citation. A source-bundled action reads one expert--publication declaration and obtains the corresponding role decisions for all applicable criteria. If the endpoints agree, further reads cannot change the certified output. Our reference implementations use deterministic action order. A learned action policy would be a separate comparison.

\subsection{Deferred numerical rechecks}
The live reviewer experiment uses the logit difference from an open cross-encoder. Dynamic tensor shapes reduce padding but can change numerical scores. Immediate rechecking recomputes every tentative pair with a fixed shape before taking another action. Deferred rechecking keeps those pairs unresolved and resolves only tentative witnesses in a reachable output. All methods use the same 0.25 logit band, checkpoint, candidate inventory and output rule within the comparison. The certificate assumes that trusted decisions outside the band have the same sign as the reference. The band is an empirical guard on the measured workload.

\subsection{Grounding and conditional transfer}
A return should distinguish a historical source fact, a candidate's role, and a proposed bridge to a target operation. A bridge identifies what input, operation and output can carry over, and which changes in equipment or context still require validation. Future project thresholds remain obligations unless a comparable prior test is documented.

For a numerical criterion, the verifier selects sentences reporting the same object, metric and conditions. It checks personal execution separately when that role is required. Exact decimal comparison converts compatible units and applies strict or inclusive inequalities to an original measurement. A design target or a colleague's result cannot establish the candidate's own performance. Conflicting comparable measurements and missing scope both yield unresolved output in this verifier. This reading stage precedes the monotone support decision in Eq.~\ref{eq:certificate}.

\section{Data and evaluation}
\label{sec:data}
\paragraph{Enterprise evaluation.}
The private pool contains 100 complete requests and 700 preselected expert views. The 80 development and 20 sealed requests have disjoint groups by expert and source identifiers in the frozen snapshot. Each request receives five dense and two random candidates, with up to three papers and one profile per expert. This pool can measure judgments conditional on a common inventory. It cannot estimate full-corpus expert recall. The main enterprise replay here uses eight development requests and leaves the sealed partition untouched. The 48-case review packet deliberately oversamples diagnostic source views and awaits independent judgments.

\paragraph{Synthesis and controlled evidence.}
The separate construction corpus contains 1,000 parent needs and 5,000 Chinese--English query pairs. Pairs remain grouped by parent and require independent semantic review before being treated as parallel ground truth. The public numerical controls hold most evidence text fixed while changing a number, a unit, an operating condition, a measurement status, or the person credited. Paired sentence-order reversals measure robustness to presentation. Publisher-statement controls join the same person across two PLOS records using fixture-local tokens for validated identifiers and vary individual CRediT requirements. These controls identify specific false-confirmation mechanisms on known inputs.

\paragraph{Public first-query panels.}
The joint nomination benchmark uses 160 PLOS requests and 51 eLife heading requests. It composes two topic--role conditions from publisher metadata. A request supplies no author or title. Each verified response contains at most one person and two citations. A known match has both intended role statements and both publisher topic labels. Missing labels leave topical support unresolved. Paired diagnostics account for repeated topic pairs. We separately replay controlled prose using held-out role descriptions. The two templates reuse the base requests and do not constitute independent replications.

\paragraph{Quality and cost.}
Report expert relevance, citation support, personal attribution and hard-constraint satisfaction against the same source views. For a fixed cutoff $K$, let $z_i=1$ when recommendation $i$ is relevant and has the required, personally attributable support. With $w_i=1/\log_2(i+1)$, the normalized gain is
\begin{equation}
G_K=\frac{\sum_{i=1}^{K} w_i z_i}{\sum_{i=1}^{K}w_i}.
\end{equation}
An unreturned position contributes zero. Supported P@5 is the unweighted mean of the five slot labels. For the stricter citation measure, every citation must substantiate the expert and criterion to which it is attached, and each required criterion must have at least one witness. A document may cover a subset of the criteria. We also report the expert-only rate, whether at least one citation supports the expert, and citation precision. An unresolved judgment contributes an interval from zero to one. In a paired contrast, identical query--expert--evidence views share the same unknown label, so identical positional weights cancel. These are identification bounds conditional on the available judgments. Sampling uncertainty calls for a separate cluster-level analysis by company, parent query, or paper. Report the rate of unsupported approvals under missing or conflicting evidence, and compare end-to-end P95 only at predeclared quality and risk targets. Record throughput and the cost per independently verified successful request alongside latency, including fallbacks.

\section{Results}
\label{sec:results}
\subsection{Joint evidence at equal source reads}
All routes open three local declarations after live 0.6B query encoding and hybrid ranking. The sequential control reads one result per condition before choosing a third source using the observed role gap. Joint nomination recovers more known matches at comparable warm serial latency but also makes more unresolved returns (Table~\ref{tab:joint}). Cached batch embeddings yield 26 eLife matches; the live handler yields 14 with unchanged downstream rules. We use the live count in comparisons. At depth 64, owner postings reproduce the exhaustive plans while lowering planning P95 from 323.29--326.09 to 109.55--110.54~$\mu$s. At depth eight, the exhaustive scan is faster.

\begin{table*}[t]
\caption{Live local search on publisher-declaration benchmarks. Each route reads three declarations and returns at most one expert per request. Known matches satisfy both topic tags and both declared roles. Unresolved refers to incomplete topic labels. P95 includes live encoding, retrieval, file reads, and serialization on one RTX 5090.}
\label{tab:joint}
\centering\small
\begin{tabular}{@{}llrrr@{}}\toprule
Panel & First-round allocation & Known matches & Unresolved returns & P95 (ms) \\\midrule
PLOS, 160 requests & Topical order & 1 & 0 & 9.68--9.76 \\
 & Sequential role-gap repair & 28 & 5 & 9.79--9.85 \\
 & Joint by contributor & 49 & 29 & 9.79--10.10 \\\midrule
eLife, 51 requests & Topical order & 0 & 0 & 10.04--10.50 \\
 & Sequential role-gap repair & 0 & 0 & 10.21--10.43 \\
 & Joint by contributor & 14 & 2 & 10.23--10.61 \\
\bottomrule
\end{tabular}
\end{table*}

A matched-issuance diagnostic within joint nomination compares rank confidence with a full topic reader. At 32 PLOS returns, the two yield 29 and 25 known matches, respectively. The paired 95\% interval for the difference per request spans $-6.54$ to $+14.12$ percentage points. Local P95 is about 10 versus 41 ms. Locked PLOS cutoffs transfer to 12 known returns for each route on eLife, while the PLOS held-out folds exceed the development target for missing topic labels. This diagnostic compares verification rules, and the quality advantage of joint nomination over sequential repair at equal issuance awaits blind review.

\subsection{Candidate depth, attributed coverage and ranking}
In the isolated PLOS replay, sequential repair and joint byline nomination both issue on 24 requests. They recover 21 and 19 known matches, respectively. Twenty outputs are identical and four change the expert. Equalizing the panel-wide count at 33 gives 28 versus 27 known matches. eLife has no shared issuing requests for this comparator. These diagnostics separate finding evidence on new queries from replacing an expert already recommended. The changed PLOS outputs have been frozen for independent audit.

Combine both catalogs without duplicate titles, DOIs or shared verified identifiers, then rerun the original 211 requests against the 452-article collection. Both publishers provide distractors to either panel. Cross-publisher results have no aligned topic annotations. Table~\ref{tab:mixed} decomposes the effect of enlarging the candidate pools. At depth 128, PLOS has 115 requests with a complete known proof in the pools. The role-aware three-read plan retains at least one such proof on 89 requests and returns one on 85. In eLife the corresponding counts are 49, 46 and 46. The deeper route yields the same expert and citations on requests both widths answer. At the shallower issuance quotas the deeper role route yields the same known counts. A fixed 32-per-publisher quota reduces proof availability in both panels.

\begin{table*}[t]
\caption{Mixed publisher retrieval with exactly three declaration reads. The method columns show known matches / issued recommendations. Reachable counts requests with at least one labeled proof available in the candidate pools.}
\label{tab:mixed}
\small\centering
\begin{tabular}{@{}llrrrr@{}}\toprule
Panel & Depth per condition & Reachable & Sequential repair & Joint byline & Joint personal roles \\\midrule
PLOS (160) & 64 & 79 & 12 / 15 & 43 / 67 & 68 / 104 \\
           & 128 & 115 & 13 / 22 & 45 / 73 & 85 / 143 \\
eLife (51) & 64 & 17 & 0 / 0 & 14 / 15 & 17 / 18 \\
           & 128 & 49 & 0 / 0 & 14 / 16 & 46 / 49 \\\bottomrule
\end{tabular}
\end{table*}

\subsection{Equal-output savings}
With complete publisher role fields, planning from depth 64 and certifying where possible returns the fixed depth-128 outcomes while reducing total reads from 633 to 516. The same response parity holds after masking or rotating indexed roles, although fewer requests qualify for early stopping. In the three 50\% exposure masks with role rotation, 15--50 PLOS and 1--4 eLife requests stop at depth 64. Warm request latency remains similar because both arms perform the deep retrieval.

In the live request control with the same padded GPU planner, byline precheck reduces reads from 1,266 to 910 per 422 requests while returning identical bytes. The source service uses six loopback read slots with 16 ms of inserted delay. At 60 scheduled arrivals/s, both policies have P95 near 33 ms. At 165 arrivals/s, all-three parallel reading backs up and yields 977--981 ms P95. The precheck yields 52--53 ms as source traffic falls. These are finite, controlled service traces. A separate source-only test finds that waiting for two declarations before issuing the third increases low-load P95 from about 18 to 33 ms versus the eager precheck.

We also compare exact reciprocal-rank fusion implementations. On the 452-article catalog, direct score arrays and vectorized fusion reduce full warm P95 by about 0.52 ms and retain every ordered list and serialized response. On 64,183 LitSearch articles, packed dense keys and a score bound preserve all 246 ordered top-128 lists. A prior vectorized partial selector takes 7.71--7.83 ms P95, packed full ordering takes 5.44--5.53 ms, and the bounded version takes 5.03--5.07 ms. This component starts after dense scoring and still calculates dense scores for all eligible documents.

\subsection{Interpretation of controlled prose}
With controlled prose, fused and separate encoding each recover 41 PLOS matches versus 49 with supplied roles. The fused path recovers both roles for 110 of 160 requests and issues five returns whose declarations lack an intended role. Warm P95 falls from 19.06--19.67 to 9.91--10.09 ms. An open 4B interpreter recovers 31 matches at 57.55--57.83 ms. Live eLife gives 10 fused matches versus 14 with supplied roles, with one fused return lacking an intended role.

\subsection{Transfer checks on reviewer search}
On a separate reviewer benchmark, a fixed-inventory certificate with deferred numerical checks reproduces exhaustive reviewer order and citations. Requiring two sources, the deferred route lowers P95 by 13.4\% on 34 CoF NIPS requests and by 13.8\% on 59 SciRepEval requests relative to the fastest fixed control. The SciRepEval paired interval includes zero. These reviewer labels judge author relevance. The supplement reports the latency, queues and evidence limits of this transfer.

\subsection{Multiple constraints and enterprise coverage}
A replay with 108 held-out publisher-statement views evaluates three role constraints across two publications. All methods reproduce exhaustive Top-1 and Top-3 order and citation IDs. For specified publications with both visible, Top-1 reads average 3.67 under exhaustive verification, 2.22 under request-order short circuit, and 1.78 under both the certificate and a fixed schedule ordered by the number of constraints touched per source. This establishes correctness on the supplied declarations. Both routes reach the minimum read count over all subsets of observed declarations on the supplied fixtures. On TREC Enterprise, which provides separate expert and source judgments, the best fixed schedule is within 1.4\% of the minimum expert--document checks for either support count under a perfect-verifier replay.

The private development replay uses eight requests, two verbatim hard requirements each, 56 candidates and 209 candidate--source views. A frozen 0.6B reranker recognizes support in one of 112 candidate--requirement slots at the 0.90 operating point. No candidate passes both requirements. At 0.50, one candidate passes both and the adaptive and strongest fixed schedule tie at 26 reads per request. The reranker score is an operational decision, not a calibrated probability or a human label. The low acceptance rate could reflect thin evidence, verifier misses, or both. Without adjudicated positives, output agreement alone cannot establish expert-search quality.

\subsection{Scope and personal attribution}
In a 480-case counterfactual control, removing the person or project check each creates 48 false approvals while leaving all intact threshold pairs correctly classified. With prose distractors, a frozen 4B selector often picks an irrelevant evidence card before arithmetic. At least 27 of 58 numerically correct decisions in that condition use a wrong card. The supplement details the bilingual measurement and attribution controls. The findings motivate separate evaluation of calculation, source scope and personal contribution.

\subsection{Shared normalization and residual obligations}
In a separate structured eLife replay, normalizing duties once per request rather than once per source lowers warm P95 from 267--270 to 85--86 ms and preserves all 72 ordered responses. With that optimization shared, role-aware nomination adds two tag-backed slots at natural issuance and gives zero gain at equal issuance. The supplement gives the latency and controls. In 36 authored requests with a mandatory condition embedded in the role clause, agreement between two role encodings still bypasses 12 conditions. Role interpretation remains a cue for source selection while the verifier retains the original obligations.

\subsection{Citation correctness at fixed expert rank}
\label{sec:citation-correctness}
We use the released expert and support judgments for all 49 TREC Enterprise topics. A shared document can support one relevant expert and provide no support for another relevant expert. We observe 1,741 such topic--document pairs across 44 topics. To isolate the metric, we construct a fixed five-person ranking with a leave-one-topic-out prior over judged candidates. We attach a known valid source and change only the second citation for the 50 slots with both a valid alternative and a judged mismatched source. The mismatched paper supports a different candidate. It has neutral support for the intended expert in 47 slots and negative support in three.

The at-least-one support rule credits all 50 slots in both conditions. The clean-citation rule credits all 50 with two supporting witnesses and zero with a mismatched second witness. These pooled counterfactuals measure metric sensitivity. Deployment error rates remain unmeasured.

We also evaluate shared-paper downweighting on live KDD and SciRepEval reviewer retrieval. Adjusting both reviewer order and evidence selection reduces repeated citation slots from 124 to 77 and from 8 to 5, respectively. Newly retrieved known positives never reach the cited top five. Holding reviewer order fixed, reweighting only citation selection changes 35 of 610 KDD and 9 of 295 SciRepEval citation packages. Even if all potentially beneficial changes were correct, they could increase clean-citation P@5 by at most 4.26 and 2.37 percentage points on these panels. This is an upper bound from the changed outputs. A blinded packet contains 118 distinct citation cases with separate content and contribution judgments pending.

\section{Limitations and publication criteria}
The public first-query panels compose tasks from publisher metadata and use verified author identifiers. Publisher topic labels are incomplete and roles may omit real contributions. Structured and templated requests leave open segmentation, personal identity under missing identifiers, enterprise qualifications, and transfer of expertise. Independent judgments of the changed public recommendations are pending. The separate sealed enterprise protocol estimates full-index versus shortlist-restricted rescue after abstention. Joint first-round nomination requires its own independent enterprise comparison. The certificate applies only to fixed inventories with monotone decisions. Public records may occur in model pretraining, and loopback HTTP leaves production network and storage costs unmeasured.

A subsequent enterprise comparison should hold interpretation, verification, and load fixed across joint nomination, sequential repair, and fixed hybrid retrieval. Match source reads and separately match recommendation counts. Blindly judge changed recommendations and their supporting spans. On independent enterprise requests, estimate supported P@5, coverage, and risk at a preregistered operating point. Evaluate client latency and cost per supported success at comparable quality, including fallback and failure.

Continual updates remain secondary. Compare them with a frozen initial agent on disjoint old and new tasks, including feedback and training costs.

\bibliographystyle{template/ACM-Reference-Format}
\bibliography{../references}
\end{document}
